\documentclass[11pt]{article}
\usepackage[margin=1in]{geometry}
\usepackage[T1]{fontenc}
\usepackage[utf8]{inputenc}
\usepackage{amsmath,amssymb,amsthm}
\usepackage{enumitem}

\title{ICPC World Finals 2020\\D. Gene Folding}
\author{}
\date{}

\begin{document}
\maketitle

\section*{Problem Summary}

We are given a string over \texttt{A}, \texttt{C}, \texttt{G}, \texttt{T}.
A fold chooses a cut between two adjacent characters such that the string is mirrored around that cut
until one end of the current string is reached.
After folding, the shorter side disappears into the longer side, so the current string becomes either a
prefix or a suffix of the old one.

We must find the minimum possible final length after any number of folds.

\section*{Key Observations}

\begin{itemize}[leftmargin=*]
    \item A fold from the left boundary of the current string is possible exactly when the current string has an
    \emph{even palindrome prefix}. If that palindrome has length $2k$, then folding removes the first $k$
    characters and keeps the suffix starting at position $k$.
    \item So, if the current left boundary is at position $\ell$ in the original string, we may jump it to a
    larger position $c$ whenever the substring
    \[
        s[\ell \dots 2c-\ell-1]
    \]
    is an even palindrome centered at $c$.
    \item Let $d_2[c]$ be the usual Manacher radius for the even palindrome centered at $c$, so the largest such
    palindrome is
    \[
        s[c-d_2[c] \dots c+d_2[c]-1].
    \]
    Then the jump $\ell \to c$ is possible iff
    \[
        c - d_2[c] \le \ell.
    \]
    \item Therefore we can scan centers from left to right and maintain the furthest left boundary
    \texttt{reach} obtainable by repeatedly folding from the left:
    whenever a center $c$ satisfies $c-d_2[c] \le \texttt{reach}$, we can extend
    \texttt{reach} to $c$.
    \item The same idea works from the right by reversing the remaining string.
    After we remove the maximal foldable prefix, the rest of the work is purely symmetric suffix-folding.
\end{itemize}

\section*{Algorithm}

\begin{enumerate}[leftmargin=*]
    \item Run Manacher's algorithm for even palindromes on the original string and compute the maximum position
    $L$ reachable by left-folds only:
    start with \texttt{reach = 0}; for every center $c$ from left to right, if
    $c-d_2[c] \le \texttt{reach}$, set \texttt{reach = c}.
    At the end, $L=\texttt{reach}$.
    \item Delete that maximal removable prefix, so the current candidate string is
    \[
        t = s[L \dots n-1].
    \]
    \item Reverse $t$, run the same left-scan again, and obtain the maximum removable prefix length $R$ of the
    reversed string. In the original orientation, this means we can remove a suffix of length $R$ from $t$.
    \item The answer is
    \[
        |t| - R.
    \]
\end{enumerate}

\section*{Correctness Proof}

We prove that the algorithm returns the correct answer.

\paragraph{Lemma 1.}
Suppose the current string is the suffix $s[\ell \dots n-1]$ of the original string.
A fold from the left boundary moves the boundary from $\ell$ to $c$ if and only if
$c-d_2[c] \le \ell$.

\paragraph{Proof.}
Such a fold is valid exactly when the prefix
$s[\ell \dots 2c-\ell-1]$ is an even palindrome centered at $c$.
This happens exactly when the maximal even palindrome centered at $c$ reaches at least to $\ell$, i.e.
\[
    c-d_2[c] \le \ell.
\]
When that condition holds, folding removes the first $c-\ell$ characters and the new left boundary is $c$.
\qed

\paragraph{Lemma 2.}
Scanning centers from left to right with the rule
``if $c-d_2[c] \le \texttt{reach}$ then set $\texttt{reach}=c$''
computes the furthest boundary reachable by any sequence of left-folds.

\paragraph{Proof.}
By Lemma 1, every successful left-fold must jump from the current boundary $\ell$ to some center $c$ with
$c-d_2[c] \le \ell$.
So whenever the scan updates \texttt{reach} to $c$, that jump is genuinely possible.
Therefore every value produced by the scan is reachable.

Conversely, consider any sequence of left-folds.
At each step its next center must satisfy $c-d_2[c] \le \ell$, where $\ell$ is the current boundary.
Since the scan always stores the greatest reachable boundary seen so far, it can make every jump that the
sequence can make, and never ends earlier.
Hence the final scan value is exactly the furthest reachable left boundary. \qed

\paragraph{Lemma 3.}
Let $L$ be the furthest left boundary computed by Lemma 2.
There exists an optimal folding process whose first phase removes exactly the prefix $s[0 \dots L-1]$.

\paragraph{Proof.}
Any left-fold only uses a mirrored prefix of the current string, so removing characters from the right never
helps a future left-fold; it can only make it harder by shortening the string.
Therefore every left-boundary movement that appears in any optimal process can already be performed before
all right-folds begin.
By Lemma 2, the furthest boundary obtainable in that left-fold phase is exactly $L$.
So some optimal process may start by removing the whole maximal foldable prefix and then continue from the
suffix $s[L \dots n-1]$. \qed

\paragraph{Lemma 4.}
After the prefix $s[0 \dots L-1]$ is removed, no further left-fold is possible, and the remaining optimum is
obtained by removing the largest possible foldable suffix.

\paragraph{Proof.}
If a left-fold were still possible on $s[L \dots n-1]$, Lemma 1 would give another valid jump beyond $L$,
contradicting the maximality from Lemma 2.
So only right-folds remain.
Reversing the string turns right-folds into left-folds, and Lemma 2 applies symmetrically.
Hence running the same scan on the reversed remainder computes exactly how many characters can still be
removed from the right. \qed

\paragraph{Theorem.}
The algorithm outputs the correct answer for every valid input.

\paragraph{Proof.}
By Lemma 2, the first scan computes the maximum prefix removable by repeated left-folds.
By Lemma 3, some optimal solution begins by removing exactly that prefix.
By Lemma 4, the rest of the optimal process is just the symmetric suffix case on the remaining string, which
the second scan computes after reversal.
Therefore the algorithm removes exactly the maximum possible prefix and suffix, and the remaining middle
segment has minimum possible length. \qed

\section*{Complexity Analysis}

Each Manacher run is linear.
We run it twice, once on the original string and once on the reversed remainder, so the total running time is
$O(n)$.

The palindrome-radius array stores $O(n)$ integers, so the memory usage is $O(n)$.

\section*{Implementation Notes}

\begin{itemize}[leftmargin=*]
    \item We use the direct even-palindrome version of Manacher, so center $c$ is the cut between
    positions $c-1$ and $c$.
    \item The scan condition is exactly \texttt{center - d2[center] <= reach}.
    \item After removing the maximal prefix, reversing the remainder lets us reuse the same code for the suffix
    phase.
\end{itemize}

\end{document}
